% Options for packages loaded elsewhere
\PassOptionsToPackage{unicode}{hyperref}
\PassOptionsToPackage{hyphens}{url}
%
\documentclass[
]{article}
\usepackage{amsmath,amssymb}
\usepackage{iftex}
\ifPDFTeX
  \usepackage[T1]{fontenc}
  \usepackage[utf8]{inputenc}
  \usepackage{textcomp} % provide euro and other symbols
\else % if luatex or xetex
  \usepackage{unicode-math} % this also loads fontspec
  \defaultfontfeatures{Scale=MatchLowercase}
  \defaultfontfeatures[\rmfamily]{Ligatures=TeX,Scale=1}
\fi
\usepackage{lmodern}
\ifPDFTeX\else
  % xetex/luatex font selection
\fi
% Use upquote if available, for straight quotes in verbatim environments
\IfFileExists{upquote.sty}{\usepackage{upquote}}{}
\IfFileExists{microtype.sty}{% use microtype if available
  \usepackage[]{microtype}
  \UseMicrotypeSet[protrusion]{basicmath} % disable protrusion for tt fonts
}{}
\makeatletter
\@ifundefined{KOMAClassName}{% if non-KOMA class
  \IfFileExists{parskip.sty}{%
    \usepackage{parskip}
  }{% else
    \setlength{\parindent}{0pt}
    \setlength{\parskip}{6pt plus 2pt minus 1pt}}
}{% if KOMA class
  \KOMAoptions{parskip=half}}
\makeatother
\usepackage{xcolor}
\usepackage[margin=1in]{geometry}
\usepackage{graphicx}
\makeatletter
\def\maxwidth{\ifdim\Gin@nat@width>\linewidth\linewidth\else\Gin@nat@width\fi}
\def\maxheight{\ifdim\Gin@nat@height>\textheight\textheight\else\Gin@nat@height\fi}
\makeatother
% Scale images if necessary, so that they will not overflow the page
% margins by default, and it is still possible to overwrite the defaults
% using explicit options in \includegraphics[width, height, ...]{}
\setkeys{Gin}{width=\maxwidth,height=\maxheight,keepaspectratio}
% Set default figure placement to htbp
\makeatletter
\def\fps@figure{htbp}
\makeatother
\setlength{\emergencystretch}{3em} % prevent overfull lines
\providecommand{\tightlist}{%
  \setlength{\itemsep}{0pt}\setlength{\parskip}{0pt}}
\setcounter{secnumdepth}{-\maxdimen} % remove section numbering
\usepackage{booktabs}
\usepackage{longtable}
\usepackage{array}
\usepackage{multirow}
\usepackage{wrapfig}
\usepackage{float}
\usepackage{colortbl}
\usepackage{tabu}
\usepackage{threeparttable}
\usepackage{threeparttablex}
\usepackage[normalem]{ulem}
\usepackage[utf8]{inputenc}
\usepackage{makecell}
\usepackage{xcolor}
\usepackage{pdflscape}
\newcommand{\blandscape}{\begin{landscape}}
\newcommand{\elandscape}{\end{landscape}}
\ifLuaTeX
  \usepackage{selnolig}  % disable illegal ligatures
\fi
\usepackage{bookmark}
\IfFileExists{xurl.sty}{\usepackage{xurl}}{} % add URL line breaks if available
\urlstyle{same}
\hypersetup{
  pdftitle={Social Signaling and Prosocial Behavior},
  pdfauthor={Anne Karing (UChicago) and Karim Naguib ()},
  hidelinks,
  pdfcreator={LaTeX via pandoc}}

\title{Social Signaling and Prosocial Behavior}
\usepackage{etoolbox}
\makeatletter
\providecommand{\subtitle}[1]{% add subtitle to \maketitle
  \apptocmd{\@title}{\par {\large #1 \par}}{}{}
}
\makeatother
\subtitle{Experimental Evidence in Community Deworming in Kenya}
\author{Anne Karing (UChicago) and Karim Naguib ()}
\date{12/10/2020}

\begin{document}
\maketitle

\begin{figure}
\includegraphics[width=\maxwidth]{/home/ed/projects/takeup/presentations/takeup-STRUCTURAL_LINEAR_U_SHOCKS_PHAT_MU_REP_FOB95-fig/nonparametric-takeup-plot-control-bracelet-1.pdf}
\caption{Nonparametric Take-up Probability}
\textit{Note:} {\footnotesize }
\end{figure}

\begin{figure}
\includegraphics[width=\maxwidth]{/home/ed/projects/takeup/presentations/takeup-STRUCTURAL_LINEAR_U_SHOCKS_PHAT_MU_REP_FOB95-fig/nonparametric-takeup-plot-control-ink-1.pdf}
\caption{Nonparametric Take-up Probability - Control and Ink}
\textit{Note:} {\footnotesize }
\end{figure}

\begin{figure}
\includegraphics[width=\maxwidth]{/home/ed/projects/takeup/presentations/takeup-STRUCTURAL_LINEAR_U_SHOCKS_PHAT_MU_REP_FOB95-fig/nonparametric-takeup-plot-control-calendar-1.pdf}
\caption{Nonparametric Take-up Probability - Control and Calendar}
\textit{Note:} {\footnotesize }
\end{figure}

\begin{figure}
\includegraphics[width=\maxwidth]{/home/ed/projects/takeup/presentations/takeup-STRUCTURAL_LINEAR_U_SHOCKS_PHAT_MU_REP_FOB95-fig/nonparametric-takeup-plot-bracelet-calendar-1.pdf}
\caption{Nonparametric Take-up Probability - Bracelet and Calendar}
\textit{Note:} {\footnotesize }
\end{figure}

\begin{figure}
\includegraphics[width=\maxwidth]{/home/ed/projects/takeup/presentations/takeup-STRUCTURAL_LINEAR_U_SHOCKS_PHAT_MU_REP_FOB95-fig/nonparametric-takeup-plot-all-1.pdf}
\caption{Nonparametric Take-up Probability - All}
\textit{Note:} {\footnotesize }
\end{figure}

\begin{figure}
\includegraphics[width=\maxwidth]{/home/ed/projects/takeup/presentations/takeup-STRUCTURAL_LINEAR_U_SHOCKS_PHAT_MU_REP_FOB95-fig/takeup-dynamics-plot-1.pdf}
\caption{Deworming Take-up By Day of Campaign}
\textit{Note:} {\footnotesize }
\end{figure}

\subsection{Component checks}\label{component-checks}

\subsection{Fit Loading}\label{fit-loading}

\paragraph{Delta Expectation}\label{delta-expectation}

\begin{figure}
\includegraphics[width=\maxwidth]{/home/ed/projects/takeup/presentations/takeup-STRUCTURAL_LINEAR_U_SHOCKS_PHAT_MU_REP_FOB95-fig/alternative-delta-with-u-example-1.pdf}
\caption{Net reputational return, $\mu\Delta[w^*]$, as a function of take-up}
\textit{Note:} {\footnotesize This plot shows how expected reputational return, $\mu\Delta[w^*]$
 varies as cutoff type, $w^*$, changes. As $\sigma_u$ increases, the
 reputational return falls as it becomes harder to infer someone's type based
 off actions.}
\end{figure}

\begin{figure}
\includegraphics[width=\maxwidth]{/home/ed/projects/takeup/presentations/takeup-STRUCTURAL_LINEAR_U_SHOCKS_PHAT_MU_REP_FOB95-fig/delta-with-u-example-1.pdf}
\caption{Net reputational return, $\mu\Delta[w^*]$, as a function of cutoff type, $w^*$}
\textit{Note:} {\footnotesize This plot shows how expected reputational return, $\mu\Delta[w^*]$
 varies as cutoff type, $w^*$, changes. As $\sigma_u$ increases, the
 reputational return falls as it becomes harder to infer someone's type based
 off actions.}
\end{figure}

\begin{verbatim}
## [1] 4
\end{verbatim}

\begin{figure}
\includegraphics[width=\maxwidth]{/home/ed/projects/takeup/presentations/takeup-STRUCTURAL_LINEAR_U_SHOCKS_PHAT_MU_REP_FOB95-fig/delta-with-u-mu-change-example-1.pdf}
\caption{Net reputational return, $\mu\Delta[w^*]$, as a function of cutoff type, $w^*$}
\textit{Note:} {\footnotesize }
\end{figure}

\begin{verbatim}
## [1] 4
\end{verbatim}

\paragraph{Reported Social Perception}\label{reported-social-perception}

\begin{figure}
\includegraphics[width=\maxwidth]{/home/ed/projects/takeup/presentations/takeup-STRUCTURAL_LINEAR_U_SHOCKS_PHAT_MU_REP_FOB95-fig/reported-social-perception-1.pdf}
\caption{Reported social perception of some observable activities.}
\textit{Note:} {\footnotesize This plot shows the praise and stigma associated with various activities
 including deworming.}
\end{figure}

\paragraph{Distance Randomisation}\label{distance-randomisation}

\begin{figure}
\includegraphics[width=\maxwidth]{/home/ed/projects/takeup/presentations/takeup-STRUCTURAL_LINEAR_U_SHOCKS_PHAT_MU_REP_FOB95-fig/dist-plot-1.pdf}
\caption{Distance to Treatment Location (meters)}
\textit{Note:} {\footnotesize This graph shows a density plot of distance to the assigned point of
 treatment from each cluster centroid and household location. Clusters assigned to
 the close condition have greater density <1.25km whilst those assigned to the
 far condition place greater density >1.25km.
 Density estimates use all 9,805 surveyed individuals in the takeup sample.}
\end{figure}

\begin{figure}
\includegraphics[width=\maxwidth]{/home/ed/projects/takeup/presentations/takeup-STRUCTURAL_LINEAR_U_SHOCKS_PHAT_MU_REP_FOB95-fig/dist-hist-1.pdf}
\caption{Distance to treatment location (meters)}
\textit{Note:} {\footnotesize This graph shows a histogram plot of distance to the assigned point of
 treatment from each cluster centroid and household location. Clusters assigned to
 the close condition have greater density <1.25km whilst those assigned to the
 far condition place greater density >1.25km.
 density estimates use all 9,805 surveyed individuals in the takeup sample.}
\end{figure}

\begin{figure}
\includegraphics[width=\maxwidth]{/home/ed/projects/takeup/presentations/takeup-STRUCTURAL_LINEAR_U_SHOCKS_PHAT_MU_REP_FOB95-fig/pot-distance-overlap-1.pdf}
\caption{Distribution of distance to closest possible point of treatment, by distance condition}
\textit{Note:} {\footnotesize This plot shows there's significant overlap between the distribution of closest points of treatment
 across distance condition - i.e. any village assigned to a far pot could have been assigned to a close pot and vice versa.}
\end{figure}

\paragraph{belief facet plots}\label{belief-facet-plots}

\begin{figure}
\includegraphics[width=\maxwidth]{/home/ed/projects/takeup/presentations/takeup-STRUCTURAL_LINEAR_U_SHOCKS_PHAT_MU_REP_FOB95-fig/fob-beliefs-line-plot-post-1.pdf}
\caption{First order belief levels and average treatment effects}
\textit{Note:} {\footnotesize This plot shows first estimates using the first order beliefs
  binomial model. Total sample size: 999 observations at the individual level.}
\end{figure}

\begin{figure}
\includegraphics[width=\maxwidth]{/home/ed/projects/takeup/presentations/takeup-STRUCTURAL_LINEAR_U_SHOCKS_PHAT_MU_REP_FOB95-fig/sob-beliefs-line-plot-post-1.pdf}
\caption{Second order belief levels and average treatment effects}
\textit{Note:} {\footnotesize This plot shows first estimates using the second order beliefs
  binomial model. total sample size: 999 observations at the individual level.}
\end{figure}

\begin{verbatim}
## [1] 2
\end{verbatim}

\paragraph{beliefs plots}\label{beliefs-plots}

\begin{figure}
\includegraphics[width=\maxwidth]{/home/ed/projects/takeup/presentations/takeup-STRUCTURAL_LINEAR_U_SHOCKS_PHAT_MU_REP_FOB95-fig/belief-2ord-ate-1.pdf}
\caption{second order belief average treatment effects}
\textit{Note:} {\footnotesize this plot shows first estimates using the second order beliefs
  binomial model. total sample size: 999 observations at the individual level.}
\end{figure}

\begin{figure}
\includegraphics[width=\maxwidth]{/home/ed/projects/takeup/presentations/takeup-STRUCTURAL_LINEAR_U_SHOCKS_PHAT_MU_REP_FOB95-fig/belief-1ord-prop-1.pdf}
\caption{first order belief levels}
\textit{Note:} {\footnotesize this plot shows first estimates using the first order beliefs
  binomial model. total sample size: 999 observations at the individual level.}
\end{figure}

\begin{figure}
\includegraphics[width=\maxwidth]{/home/ed/projects/takeup/presentations/takeup-STRUCTURAL_LINEAR_U_SHOCKS_PHAT_MU_REP_FOB95-fig/belief-2ord-prop-1.pdf}
\caption{second order belief levels}
\textit{Note:} {\footnotesize this plot shows first estimates using the second order beliefs
  binomial model. total sample size: 999 observations at the individual level.}
\end{figure}

\subsection{RF Belief Model}\label{rf-belief-model}

\begin{figure}
\includegraphics[width=\maxwidth]{/home/ed/projects/takeup/presentations/takeup-STRUCTURAL_LINEAR_U_SHOCKS_PHAT_MU_REP_FOB95-fig/fob-beliefs-line-plot-rf-1.pdf}
\caption{First order belief levels and average treatment effects}
\textit{Note:} {\footnotesize This plot shows estimates using the first order beliefs
  linear probability model. Total sample size: 999 observations at the individual level.}
\end{figure}

\begin{figure}
\includegraphics[width=\maxwidth]{/home/ed/projects/takeup/presentations/takeup-STRUCTURAL_LINEAR_U_SHOCKS_PHAT_MU_REP_FOB95-fig/sob-beliefs-line-plot-rf-1.pdf}
\caption{Second order belief levels and average treatment effects}
\textit{Note:} {\footnotesize This plot shows estimates using the second order beliefs
  linear probability model. Total sample size: 999 observations at the individual level.}
\end{figure}

\paragraph{disagg beliefs}\label{disagg-beliefs}

\begin{verbatim}
## [1] 2
\end{verbatim}

\begin{figure}
\includegraphics[width=\maxwidth]{/home/ed/projects/takeup/presentations/takeup-STRUCTURAL_LINEAR_U_SHOCKS_PHAT_MU_REP_FOB95-fig/disagg-fob-close-1.pdf}
\caption{disaggregated first order belief levels - close}
\textit{Note:} {\footnotesize this plot shows estimates using the first-order beliefs
  simple means.
  total sample size: 999 observations at the individual level.}
\end{figure}

\begin{verbatim}
## [1] 2
\end{verbatim}

\begin{figure}
\includegraphics[width=\maxwidth]{/home/ed/projects/takeup/presentations/takeup-STRUCTURAL_LINEAR_U_SHOCKS_PHAT_MU_REP_FOB95-fig/disagg-fob-far-1.pdf}
\caption{disaggregated first order belief levels - close}
\textit{Note:} {\footnotesize this plot shows estimates using the first-order beliefs
  simple means.
  total sample size: 999 observations at the individual level.}
\end{figure}

\begin{verbatim}
## [1] 2
\end{verbatim}

\begin{figure}
\includegraphics[width=\maxwidth]{/home/ed/projects/takeup/presentations/takeup-STRUCTURAL_LINEAR_U_SHOCKS_PHAT_MU_REP_FOB95-fig/disagg-fob-far-single-ci-1.pdf}
\caption{disaggregated first order belief levels - close}
\textit{Note:} {\footnotesize this plot shows estimates using the first-order beliefs
  simple means.
  total sample size: 999 observations at the individual level.}
\end{figure}

\begin{verbatim}
## [1] 2
\end{verbatim}

\begin{figure}
\includegraphics[width=\maxwidth]{/home/ed/projects/takeup/presentations/takeup-STRUCTURAL_LINEAR_U_SHOCKS_PHAT_MU_REP_FOB95-fig/disagg-fob-both-1.pdf}
\caption{disaggregated first order belief levels}
\textit{Note:} {\footnotesize this plot shows estimates using the first-order beliefs
  simple means.
  total sample size: 999 observations at the individual level.}
\end{figure}

\begin{verbatim}
## [1] 4
\end{verbatim}

\begin{figure}
\includegraphics[width=\maxwidth]{/home/ed/projects/takeup/presentations/takeup-STRUCTURAL_LINEAR_U_SHOCKS_PHAT_MU_REP_FOB95-fig/disagg-sob-both-1.pdf}
\caption{disaggregated second order belief levels}
\textit{Note:} {\footnotesize this plot shows estimates using the first-order beliefs
  simple means.
  total sample size: 999 observations at the individual level.}
\end{figure}

\begin{figure}
\includegraphics[width=\maxwidth]{/home/ed/projects/takeup/presentations/takeup-STRUCTURAL_LINEAR_U_SHOCKS_PHAT_MU_REP_FOB95-fig/pref-gift-plot-1.pdf}
\caption{preferences for gift by treatment arm}
\textit{Note:} {\footnotesize this plot shows the preference for either the bracelet or calendar gift
 by each treatment arm/distance condition.}
\end{figure}

\begin{figure}
\includegraphics[width=\maxwidth]{/home/ed/projects/takeup/presentations/takeup-STRUCTURAL_LINEAR_U_SHOCKS_PHAT_MU_REP_FOB95-fig/pref-gift-plot-2.pdf}
\caption{preferences for gift by treatment arm}
\textit{Note:} {\footnotesize this plot shows the preference for either the bracelet or calendar gift
 by each treatment arm/distance condition.}
\end{figure}

\paragraph{incentive ates}\label{incentive-ates}

\begin{verbatim}
## [1] 2
\end{verbatim}

\begin{figure}
\includegraphics[width=\maxwidth]{/home/ed/projects/takeup/presentations/takeup-STRUCTURAL_LINEAR_U_SHOCKS_PHAT_MU_REP_FOB95-fig/reduced-incentive-level-plot-1.pdf}
\caption{incentive takeup levels, reduced form model}
\textit{Note:} {\footnotesize this plot shows the level of takeup in each treatment condition by
 distance, estimated using a bayesian probit ('reduced form') model. the prior
 predictive distribution is overlaid at the top of the plot, which shows the
 predicted treatment effects before conditioning on the data. the prior is
 centered at 50\% takeup and the 95\% credible interval covers almost the entire
 unit interval. the plot shows that deworming takeup is higher in the close condition
 and that the bracelet has the highest level amongst the offered treatments.
 line range: 95\% credible interval.
 outer box: 80\% credible interval. inner box: 50\% credible interval.
 point: mean. 9,805 observations from the individual takeup sample used for estimates.}
\end{figure}

\begin{figure}
\includegraphics[width=\maxwidth]{/home/ed/projects/takeup/presentations/takeup-STRUCTURAL_LINEAR_U_SHOCKS_PHAT_MU_REP_FOB95-fig/reduced-incentive-ate-plot-1.pdf}
\caption{incentive average treatment effect, reduced form model}
\textit{Note:} {\footnotesize line range: 95\% credible interval.
 outer box: 80\% credible interval. inner box: 50\% credible interval.
 thick vertical line: median. point: mean.}
\end{figure}

\begin{verbatim}
## [1] 3
\end{verbatim}

\begin{figure}
\includegraphics[width=\maxwidth]{/home/ed/projects/takeup/presentations/takeup-STRUCTURAL_LINEAR_U_SHOCKS_PHAT_MU_REP_FOB95-fig/structural-signalling-ate-plot-1.pdf}
\caption{signalling average treatment effect, structural model}
\textit{Note:} {\footnotesize this plot shows the average signalling treatment effect estimated
 using the structural model. the prior predictive distribution is overlaid at
 the top of the plot, which shows the
 predicted signalling treatment effects before conditioning on the data. the prior is
 centered at a treatment effect of approximately 0. the plot shows that the signalling effect
 is greatest for the ink and bracelet treatments and largest in the far condition.
 line range: 95\% credible interval.
 outer box: 80\% credible interval. inner box: 50\% credible interval.
 thick vertical line: median. point: mean. signalling effect calculated using
 structural model incorporating willingness-to-pay, beliefs, and takeup data.}
\end{figure}

\begin{verbatim}
## [1] 2
\end{verbatim}

\begin{figure}
\includegraphics[width=\maxwidth]{/home/ed/projects/takeup/presentations/takeup-STRUCTURAL_LINEAR_U_SHOCKS_PHAT_MU_REP_FOB95-fig/structural-private-ate-plot-1.pdf}
\caption{private incentive average treatment effect, structural model}
\textit{Note:} {\footnotesize this plot shows the average private incentive treatment effect estimated
 using the structural model. the prior predictive distribution is overlaid at
 the top of the plot, which shows the
 predicted priavte incentive treatment effects before conditioning on the data. the prior is
 centered at a treatment effect of approximately 0. the plot shows that the private incentive effect
 is greatest for the calendar and bracelet treatments.
 line range: 95\% credible interval.
 outer box: 80\% credible interval. inner box: 50\% credible interval.
 thick vertical line: median. point: mean.}
\end{figure}

\begin{verbatim}
## [1] 5
\end{verbatim}

\begin{figure}
\includegraphics[width=\maxwidth]{/home/ed/projects/takeup/presentations/takeup-STRUCTURAL_LINEAR_U_SHOCKS_PHAT_MU_REP_FOB95-fig/structural-incentive-ate-plot-1.pdf}
\caption{incentive average treatment effect, structural model}
\textit{Note:} {\footnotesize this plot shows the average treatment effect estimated
 using the structural model. the prior predictive distribution is overlaid at
 the top of the plot, which shows the
 predicted treatment effects before conditioning on the data. the prior is
 centered at a treatment effect of 0 and the 95\% credible interval covers the
 interval $[-0.125,0.125]$. the plot shows that the treatment effect is greatest
 for the bracelet treatment and estimated treatment effects are larger for the
 bracelet at close than far.
 line range: 95\% credible interval.
 outer box: 80\% credible interval. inner box: 50\% credible interval.
 thick vertical line: median. point: mean. treatment effect calculated using
 structural model incorporating willingness-to-pay, beliefs, and takeup data.}
\end{figure}

\begin{verbatim}
## [1] 2
\end{verbatim}

\begin{figure}
\includegraphics[width=\maxwidth]{/home/ed/projects/takeup/presentations/takeup-STRUCTURAL_LINEAR_U_SHOCKS_PHAT_MU_REP_FOB95-fig/wtp-post-joint-1.pdf}
\caption{willingness-to-pay estimates}
\textit{Note:} {\footnotesize calendar valued at 47.2ksh more than bracelet with a 90\% credibility
 interval of $[41.2, 52.8]$.
 line range: 90\% credible interval.
 outer box: 80\% credible interval. inner box: 50\% credible interval.
 thick vertical line: median. point: mean.  preferences estimated using 998 individuals
 with a bayesian multinomial discrete choice model using a probit likelihood.}
\end{figure}

\paragraph{distance prob plots}\label{distance-prob-plots}

\begin{figure}
\includegraphics[width=\maxwidth]{/home/ed/projects/takeup/presentations/takeup-STRUCTURAL_LINEAR_U_SHOCKS_PHAT_MU_REP_FOB95-fig/take-up-prob-with-obs-takeup-and-linerange-1.pdf}
\caption{Structural Model Estimates of Take-up Probability by Distance}
\textit{Note:} {\footnotesize This plot shows estimated take-up probability across bracelet and
 control over distance. Take-up probability is calculated using the structural
 model's estimates of $w^*$ in each cluster and taking the complement of the Gaussian CDF
 evalutaed at $w^*$.  Outer Ribbon: 90\% credibility interval,
 inner ribbon: 50\% credibility interval.}
\end{figure}

\subsection{rep return by distance}\label{rep-return-by-distance}

\begin{figure}
\includegraphics[width=\maxwidth]{/home/ed/projects/takeup/presentations/takeup-STRUCTURAL_LINEAR_U_SHOCKS_PHAT_MU_REP_FOB95-fig/rep-return-te-dist-bracelet-1.pdf}
\caption{difference in value of reputational return across bracelet and control}
\textit{Note:} {\footnotesize value priced in terms of walking distance. estimates generated using
 decomposition from structural model.}
\end{figure}

\begin{figure}
\includegraphics[width=\maxwidth]{/home/ed/projects/takeup/presentations/takeup-STRUCTURAL_LINEAR_U_SHOCKS_PHAT_MU_REP_FOB95-fig/rep-return-te-dist-calendar-1.pdf}
\caption{difference in value of reputational return across calendar and control}
\textit{Note:} {\footnotesize value priced in terms of walking distance. estimates generated using
 decomposition from structural model.}
\end{figure}

\begin{figure}
\includegraphics[width=\maxwidth]{/home/ed/projects/takeup/presentations/takeup-STRUCTURAL_LINEAR_U_SHOCKS_PHAT_MU_REP_FOB95-fig/rep-return-te-dist-ink-1.pdf}
\caption{difference in value of reputational return across ink and control}
\textit{Note:} {\footnotesize value priced in terms of walking distance. estimates generated using
 decomposition from structural model.}
\end{figure}

\begin{figure}
\includegraphics[width=\maxwidth]{/home/ed/projects/takeup/presentations/takeup-STRUCTURAL_LINEAR_U_SHOCKS_PHAT_MU_REP_FOB95-fig/rep-return-dist-value-1.pdf}
\caption{value of reputational return across bracelet and control}
\textit{Note:} {\footnotesize value priced in terms of walking distance. estimates generated using
 decomposition from structural model.}
\end{figure}

\subsection{social multiplier}\label{social-multiplier}

\begin{figure}
\includegraphics[width=\maxwidth]{/home/ed/projects/takeup/presentations/takeup-STRUCTURAL_LINEAR_U_SHOCKS_PHAT_MU_REP_FOB95-fig/social-multiplier-1.pdf}
\caption{Estimated Social Multiplier Over Distance}
\textit{Note:} {\footnotesize this plot shows the social multiplier estimated using the structural
 model. the plot shows how sensitive take-up is to changes in cost. for the ink and
 bracelet treatment arm, as distance
 increases individuals become less and less sensitive to the increased cost due to
 greater associated reputational returns. however, for the calendar and control treatment arms
 individuals become increasingly sensitive to changes in costs.
 line: median.}
\end{figure}

\begin{figure}
\includegraphics[width=\maxwidth]{/home/ed/projects/takeup/presentations/takeup-STRUCTURAL_LINEAR_U_SHOCKS_PHAT_MU_REP_FOB95-fig/social-multiplier-uncertainty-1.pdf}
\caption{estimated social multiplier over distance}
\textit{Note:} {\footnotesize this plot shows the social multiplier estimated using the structural
 model. the plot shows how sensitive take-up is to changes in cost. for the ink and
 bracelet treatment arm, as distance
 increases individuals become less and less sensitive to the increased cost due to
 greater associated reputational returns. however, for the calendar and control treatment arms
 individuals become increasingly sensitive to changes in costs.
 line: median.
 outer ribbon: 90\% credible interval.
 inner ribbon: 50\% credible interval.}
\end{figure}

\begin{figure}
\includegraphics[width=\maxwidth]{/home/ed/projects/takeup/presentations/takeup-STRUCTURAL_LINEAR_U_SHOCKS_PHAT_MU_REP_FOB95-fig/sm-decomp-1.pdf}
\caption{Estimated Social Multiplier Over Distance - Components}
\textit{Note:} {\footnotesize This plot shows the social multiplier components estimated using the structural
 model. The plot decomposes how much of the social multiplier is driven by the
 $\Delta$ component versus how much is driven by changes in visibility
 across distance. The full line shows the entire social multiplier whilst the
 dashed line shows the $\Delta$ component - the difference between the two
 lines is driven by changes in $\mu$ visibility.
 Line: median.}
\end{figure}

\begin{figure}
\includegraphics[width=\maxwidth]{/home/ed/projects/takeup/presentations/takeup-STRUCTURAL_LINEAR_U_SHOCKS_PHAT_MU_REP_FOB95-fig/sm-decomp-annotated-1.pdf}
\caption{Estimated Social Multiplier Over Distance - Components}
\textit{Note:} {\footnotesize This plot shows the social multiplier components estimated using the structural
 model. The plot decomposes how much of the social multiplier is driven by the
 $\Delta$ component versus how much is driven by changes in visibility
 across distance. The full line shows the entire social multiplier whilst the
 dashed line shows the $\Delta$ component - the difference between the two
 lines is driven by changes in $\mu$ visibility.
 Line: median.}
\end{figure}

\begin{figure}
\includegraphics[width=\maxwidth]{/home/ed/projects/takeup/presentations/takeup-STRUCTURAL_LINEAR_U_SHOCKS_PHAT_MU_REP_FOB95-fig/rate-of-change-prior-1.pdf}
\caption{Rate of Change Estimates}
\textit{Note:} {\footnotesize This plot shows the rate of change estimated using the structural
 model. The plot shows how sensitive take-up is to changes in cost as cost (distance)
 changes. The dashed line indicates the rate of change with no visibility. Estimates
 above the dashed line show that visibility mitigates the increased cost whilst estimates
 below the dashed line indicate an amplification effect. The dashed line is
 curved due to the Probit functional form. The grey line indicates the prior
 predictive distribution - this shows what our prior beliefs imply about the
 rate of change posterior distribution before conditioning on the data.
 Line: median.
 Outer ribbon: 90\% credible interval.
 Inner ribbon: 50\% credible interval.}
\end{figure}

\begin{figure}
\includegraphics[width=\maxwidth]{/home/ed/projects/takeup/presentations/takeup-STRUCTURAL_LINEAR_U_SHOCKS_PHAT_MU_REP_FOB95-fig/rate-of-change-prior-2.pdf}
\caption{Rate of Change Estimates}
\textit{Note:} {\footnotesize This plot shows the rate of change estimated using the structural
 model. The plot shows how sensitive take-up is to changes in cost as cost (distance)
 changes. The dashed line indicates the rate of change with no visibility. Estimates
 above the dashed line show that visibility mitigates the increased cost whilst estimates
 below the dashed line indicate an amplification effect. The dashed line is
 curved due to the Probit functional form. The grey line indicates the prior
 predictive distribution - this shows what our prior beliefs imply about the
 rate of change posterior distribution before conditioning on the data.
 Line: median.
 Outer ribbon: 90\% credible interval.
 Inner ribbon: 50\% credible interval.}
\end{figure}

\begin{figure}
\includegraphics[width=\maxwidth]{/home/ed/projects/takeup/presentations/takeup-STRUCTURAL_LINEAR_U_SHOCKS_PHAT_MU_REP_FOB95-fig/rate-of-change-prior-3.pdf}
\caption{Rate of Change Estimates}
\textit{Note:} {\footnotesize This plot shows the rate of change estimated using the structural
 model. The plot shows how sensitive take-up is to changes in cost as cost (distance)
 changes. The dashed line indicates the rate of change with no visibility. Estimates
 above the dashed line show that visibility mitigates the increased cost whilst estimates
 below the dashed line indicate an amplification effect. The dashed line is
 curved due to the Probit functional form. The grey line indicates the prior
 predictive distribution - this shows what our prior beliefs imply about the
 rate of change posterior distribution before conditioning on the data.
 Line: median.
 Outer ribbon: 90\% credible interval.
 Inner ribbon: 50\% credible interval.}
\end{figure}

\begin{figure}
\includegraphics[width=\maxwidth]{/home/ed/projects/takeup/presentations/takeup-STRUCTURAL_LINEAR_U_SHOCKS_PHAT_MU_REP_FOB95-fig/rate-of-change-prior-4.pdf}
\caption{Rate of Change Estimates}
\textit{Note:} {\footnotesize This plot shows the rate of change estimated using the structural
 model. The plot shows how sensitive take-up is to changes in cost as cost (distance)
 changes. The dashed line indicates the rate of change with no visibility. Estimates
 above the dashed line show that visibility mitigates the increased cost whilst estimates
 below the dashed line indicate an amplification effect. The dashed line is
 curved due to the Probit functional form. The grey line indicates the prior
 predictive distribution - this shows what our prior beliefs imply about the
 rate of change posterior distribution before conditioning on the data.
 Line: median.
 Outer ribbon: 90\% credible interval.
 Inner ribbon: 50\% credible interval.}
\end{figure}

\subsection{sms}\label{sms}

\begin{figure}
\includegraphics[width=\maxwidth]{/home/ed/projects/takeup/temp-data/sms-noreminder/TE-by-sms-incentive.pdf}
\caption{}
\textit{Note:} {\footnotesize }
\end{figure}

\begin{figure}
\includegraphics[width=\maxwidth]{/home/ed/projects/takeup/temp-data/sms-noreminder/TE-diff-by-sms-incentive.pdf}
\caption{}
\textit{Note:} {\footnotesize }
\end{figure}

\paragraph{distance to closest school}\label{distance-to-closest-school}

\begin{figure}
\includegraphics[width=\maxwidth]{/home/ed/projects/takeup/presentations/takeup-STRUCTURAL_LINEAR_U_SHOCKS_PHAT_MU_REP_FOB95-fig/unnamed-chunk-56-1.pdf}
\caption{}
\textit{Note:} {\footnotesize }
\end{figure}

\begin{verbatim}
## # A tibble: 1 x 1
##   `n_distinct(KEY.individ)`
##                       <int>
## 1                      1512
\end{verbatim}

\begin{verbatim}
## # A tibble: 8 x 3
##   gpt_name      clean_name                   n
##   <chr>         <fct>                    <int>
## 1 Don't Know    Don't Know                1538
## 2 campaign      Direct Observation        1275
## 3 circumstances Personal Characteristics   575
## 4 communication Communication              905
## 5 other         Other                       92
## 6 relationship  Social Proximity          1425
## 7 signal        Incentive                  315
## 8 type          Personal Characteristics   333
\end{verbatim}

\begin{verbatim}
## [1] "#7FC97F" "#BEAED4" "#FDC086" "#FFFF99" "#386CB0" "#F0027F" "#BF5B17"
## [8] "#666666"
\end{verbatim}

\begin{figure}
\includegraphics[width=\maxwidth]{/home/ed/projects/takeup/presentations/takeup-STRUCTURAL_LINEAR_U_SHOCKS_PHAT_MU_REP_FOB95-fig/gpt-reason-plot-1.pdf}
\caption{Reasons Why Participants Believed Others Knew Treatment Status}
\textit{Note:} {\footnotesize }
\end{figure}

\subsubsection{Frequentist Reduced Form
Plots}\label{frequentist-reduced-form-plots}

\paragraph{ATEs}\label{ates}

\begin{figure}
\includegraphics[width=\maxwidth]{/home/ed/projects/takeup/presentations/takeup-STRUCTURAL_LINEAR_U_SHOCKS_PHAT_MU_REP_FOB95-fig/freq-rf-te-plot-1.pdf}
\caption{Incentive Takeup ATEs, Reduced Form Model}
\textit{Note:} {\footnotesize }
\end{figure}

\begin{verbatim}
## [1] 2
\end{verbatim}

\begin{figure}
\includegraphics[width=\maxwidth]{/home/ed/projects/takeup/presentations/takeup-STRUCTURAL_LINEAR_U_SHOCKS_PHAT_MU_REP_FOB95-fig/freq-rf-levels-plot-1.pdf}
\caption{Incentive Takeup Levels, Reduced Form Model}
\textit{Note:} {\footnotesize }
\end{figure}

\begin{figure}
\includegraphics[width=\maxwidth]{/home/ed/projects/takeup/presentations/takeup-STRUCTURAL_LINEAR_U_SHOCKS_PHAT_MU_REP_FOB95-fig/freq-rf-levels-discrete-dist-lasso-1.pdf}
\caption{Incentive Takeup Levels, Reduced Form Model - Discrete Distance w/ LASSO Covs}
\textit{Note:} {\footnotesize }
\end{figure}

\begin{figure}
\includegraphics[width=\maxwidth]{/home/ed/projects/takeup/presentations/takeup-STRUCTURAL_LINEAR_U_SHOCKS_PHAT_MU_REP_FOB95-fig/freq-rf-ates-discrete-dist-lasso-1.pdf}
\caption{Incentive Takeup Levels, Reduced Form Model - Discrete Distance w/ LASSO Covs}
\textit{Note:} {\footnotesize }
\end{figure}

\begin{figure}
\includegraphics[width=\maxwidth]{/home/ed/projects/takeup/presentations/takeup-STRUCTURAL_LINEAR_U_SHOCKS_PHAT_MU_REP_FOB95-fig/freq-rf-ates-recent-1.pdf}
\caption{Incentive Takeup ATEs, Reduced Form Model - Recentered}
\textit{Note:} {\footnotesize }
\end{figure}

\begin{figure}
\includegraphics[width=\maxwidth]{/home/ed/projects/takeup/presentations/takeup-STRUCTURAL_LINEAR_U_SHOCKS_PHAT_MU_REP_FOB95-fig/freq-rf-tes-discrete-plot-1.pdf}
\caption{Incentive Takeup ATEs, Reduced Form Model - Discrete Distance}
\textit{Note:} {\footnotesize }
\end{figure}

\begin{figure}
\includegraphics[width=\maxwidth]{/home/ed/projects/takeup/presentations/takeup-STRUCTURAL_LINEAR_U_SHOCKS_PHAT_MU_REP_FOB95-fig/freq-rf-tes-cts-lasso-plot-1.pdf}
\caption{Incentive Takeup ATEs, Reduced Form Model - Continuous Distance, LASSO Covs}
\textit{Note:} {\footnotesize }
\end{figure}

\paragraph{FOBs}\label{fobs}

\begin{figure}
\includegraphics[width=\maxwidth]{/home/ed/projects/takeup/presentations/takeup-STRUCTURAL_LINEAR_U_SHOCKS_PHAT_MU_REP_FOB95-fig/fob-reduced-form-1.pdf}
\caption{First Order Beliefs, Reduced Form Model}
\textit{Note:} {\footnotesize }
\end{figure}

\end{document}
